\chapter{Эквивариантные монеты физического пакета}
\label{ch:v6-spectral-transition-discrete-equivariant-coin-selector-gate}

\section{Вопрос после скалярного запрета}

Полный $M_{20}$--$M_{15}$-бимодуль разрешает только скалярную внутреннюю
монету и сохраняет кратность $300$. Однако матричные факторы были
контрольными углами, а физическая алгебра меньше. Проверим, выбирает ли её
коммутант единственную нескалярную монету.

Антикруговая сверка опирается на уже доказанные результаты Тома V:
\begin{enumerate}
 \item физический пакет одной семейной линии равен
 $H_{15}=Q_L^{(6)}\oplus L_L^{(2)}\oplus u_R^{(3)}\oplus
 d_R^{(3)}\oplus e_R^{(1)}$;
 \item коммутант сохраняющей это разложение наблюдаемой алгебры имеет
 комплексную размерность пять;
 \item голономный проектор сокращает семейную кратность, но не выбирает
 наблюдаемый блок;
 \item постоянного слабого нейтринного проектора ранга один не существует.
\end{enumerate}
Тем самым новый вопрос относится именно к локальной монете, а не повторяет
поиск семейного или хиггсовского проектора.

\section{Полная классификация секторной монеты}

Обозначим пять центральных проекторов через
\begin{equation}
 P_Q,P_L,P_u,P_d,P_e,
 \qquad
 \rank(P_s)=(6,2,3,3,1).
\end{equation}
Для максимальной наблюдаемой блочной алгебры
\begin{equation}
 \mathcal B_{\rm obs}=M_6\oplus M_2\oplus M_3\oplus M_3\oplus\mathbb C
\end{equation}
имеем
\begin{equation}
 \mathcal B_{\rm obs}'=
 \bigoplus_{s\in\{Q,L,u,d,e\}}\mathbb C P_s.
 \label{eq:v6-discrete-coin-observed-commutant}
\end{equation}
Так как физическая алгебра Стандартной модели содержится в
$\mathcal B_{\rm obs}$, её коммутант не меньше этого пространства.

На направленном носителе $\mathbb C^2_{\rm dir}\otimes H_{15}$ всякая
блочно эквивариантная монета имеет вид
\begin{equation}
 C=\sum_s C_s\otimes P_s,
 \qquad C_s\in U(2).
 \label{eq:v6-discrete-sector-coin}
\end{equation}
Следовательно, уже нижняя граница пространства монет равна пяти независимым
$U(2)$-блокам. Для однопараметрического вращения предыдущего гейта
\begin{equation}
 C(\Psi)=\sum_s
 \exp\!\left[-i\theta_s(\Psi)\sigma_y\right]\otimes P_s
 \label{eq:v6-discrete-five-angle-coin}
\end{equation}
остаётся пять углов, или четыре относительных угла после удаления общего.
Алгебра разрешает их, но не выбирает ненулевую точку этого пространства.

\section{Единственный ранг один и его физический смысл}

Среди центральных проекторов действительно имеется единственный проектор
ранга один: $P_e$. Но он выделяет правый электронный синглет $e_R$, а не
нейтральную частицу и не полный массивный электрон.

Массовое замыкание заряженного лептона связывает $e_R$ со слабым дублетом
$L_L$. До выбора Хиггса минимальная калибровочно замкнутая опора равна
\begin{equation}
 (P_L+P_e)H_{15},\qquad \rank(P_L+P_e)=3.
 \label{eq:v6-discrete-lepton-closed-rank}
\end{equation}
Выделение одной левой компоненты требует нормированного направления $H$ и
возвращает хиггс-зависимый проектор. Для нейтрино ситуация ещё строже:
коммутант полного слабого действия на $L_L\simeq\mathbb C^2$ скалярен,
поэтому постоянного эквивариантного проектора ранга один нет.

Таким образом, критерий ``выбрать единственный минимальный центральный
блок'' выбирает только хиральный конец известного юкавского ребра. Чтобы
превратить его в локализованную физическую частицу, необходимо вернуть
Хиггс и относительный коэффициент юкавской монеты. Это именно остаточная
свобода $u,d,e$, уже найденная в физическом модуле одноформ Тома V.

\begin{theorem}[Физическая редукция не даёт единственной монеты]
После ограничения полного фактора до физического пакета нескалярные
эквивариантные монеты существуют. Однако их пространство содержит не менее
пяти независимых секторных блоков и четыре относительных однопараметрических
угла. Единственный центральный проектор ранга один равен $P_e$, но он не
замкнут относительно калибровочно-юкавской динамики. Замыкание требует
$P_L+P_e$ и поля Хиггса. Поэтому физическая алгебра не выводит единственный
локальный endpoint и не снимает селекторную неоднозначность.
\end{theorem}

\begin{nogocriterion}
Нельзя объявлять $P_e$ частицей только потому, что это единственный
одномерный центральный блок. Он является правым хиральным представлением.
Добавление к монете выбранного направления внутри $L_L$ без динамического
Хиггса нарушает слабую ковариантность; добавление Хиггса возвращает уже
исследованную операторнозначную ветвь Тома V.
\end{nogocriterion}

\section{Следующий гейт}

Следующий гейт
\path{version6_spectral_transition_discrete_chiral_coin_closure_gate}
должен проверить минимальную монету на калибровочно замкнутой опоре
$L_L\oplus e_R$ и определить, выводится ли её хиральное смешивание из
локального дискретного правила либо неизбежно содержит независимые $H$ и
$y_e$.

Нулевая гипотеза: ковариантное смешивание имеет вид известного юкавского
интертвинера, зависит от Хиггса и свободной амплитуды. Стоп-критерий: если
дискретность фиксирует только шаг эволюции, но не $H$, $y_e$ и физический
масштаб, ветвь является дискретизацией старого $H_{15}$-родителя и закрывается.

\begin{protocol}
\begin{enumerate}
 \item нескалярные физически эквивариантные монеты: \textbf{существуют};
 \item число обязательных секторных блоков: \textbf{не менее пяти};
 \item относительная размерность однопараметрической монеты:
 \textbf{четыре};
 \item единственный центральный ранг один: \textbf{$e_R$};
 \item полный физический массивный endpoint из $e_R$: \textbf{нет};
 \item постоянный нейтринный ранг один: \textbf{запрещён};
 \item хиггс-зависимый селектор: \textbf{условно существует};
 \item единственная монета: \textbf{не выведена};
 \item R4 и R5: \textbf{не закрыты}.
\end{enumerate}
\end{protocol}

Машинный сертификат сохранён в
\path{s2t_v6_spectral_transition_discrete_equivariant_coin_selector_gate_results.json}.